\mysec{Lista 6}
\etocsetnexttocdepth{subsection}
\localtableofcontents
\newcounter{seis}
\stepcounter{seis}

\myexe{1}
\bigline
\textbf{Sejam $(X, \mcal F)$ um espaço mensurável e $\nu$ uma medida com sinal. Dada 
a decomposição de Jordan $\nu = \nu^+ - \nu^-$, mostre que se $\nu = \mu -\lambda$, em que 
$\mu$ e $\lambda$ são medidas positivas, com ao menos uma delas finita, então $\mu(E) \ge \nu^+(E)$
e $\lambda(E) \ge \nu^-(E)$, para todo $E \in \mcal F$.}
\begin{proof}[\textbf{Dem:}] Considere $P, N \in \mcal F$ como na decomposição de Jordan.
Suponha, sem perda de generalidade, que $\lambda(E) < +\infty$ 
para todo $E \in \mcal F$. Então, dado $E \in \mcal F$,
\begin{align*}
  \mu(E) = \nu(E) + \lambda(E) \ge \nu(E) \ge \nu(E\cap P) = \nu^+(E),
\end{align*}
Com isso, segue a primeira desigualdade. Mais ainda, note que, como $\lambda(E) < +\infty$, 
$\nu(E)$ não assume o valor $-\infty$ para nenhum $E \in \mcal F$. Desse modo, 
já que $\nu^+$ é uma medida positiva, temos $\nu^-(E) < +\infty$ para todo $E \in \mcal F$.
Desse modo, $\nu^+(E) = \nu(E) + \nu^-(E)$ para todo $E \in \mcal F$. Da desigualdade já 
obtida, 
\begin{align*}
    \mu(E) \ge \nu^+(E) \implies \nu(E)+\lambda(E) \ge \nu(E) +\nu^-(E).
\end{align*}
Daqui, se $\nu(E)< +\infty$, segue diretamente que $\lambda(E) \ge \nu^-(E)$. Por outro lado, 
se $\nu(E) = +\infty$,
\end{proof}

%ex 2 
\nextexe{seis}
\textbf{}

%ex 3 
\nextexe{seis}
\textbf{Sejam $(X, \mcal F)$ um espaço mensurável e $v$ uma medida com sinal. 
Dada a decomposição de Jordan $\nu = \nu^+ - \nu^-$, mostre que 
\begin{enumerate}[label=(\alph*)]
  \item $\nu^+(E) = \sup\set{\nu(F) \st F \subseteq E, F \in \mcal F}$. 
    \begin{proof}[\textbf{Dem:}] Para cada $E \in \mcal F$, considere
      $A_E = \set{\nu(F) \st F\subseteq E, F \in \mcal F}$. Note que, 
      dada para $P, N \in \mcal F$ como nas hipóteses da decomposição 
      de Jordan,
      \begin{align*}
        \nu^+(E) = \nu(E\cap P) \in A_E \implies \nu^+(E) \le \sup A_E.
      \end{align*}
      Por outro lado, para todo $\nu(F) \in A_E$,
      \begin{align*}
        \nu(F) = \nu^+(F) - \nu^-(F) \le \nu^+(F) \le \nu^+(E),
      \end{align*}
      desse modo, $\sup A_E \le \nu^+ (E)$, ou seja, $\sup A_E = \nu^+(E)$.
    \end{proof}
  \item $\nu^-(E) = -\inf \set{\nu(F) \st F \subseteq E, F \in \mcal F}$.
    \begin{proof}[\textbf{Dem:}] Semelhante ao item anterior,
      \begin{align*}
        -\nu^-(E) = \nu(E\cap N) \in A_E \implies -\nu^-(E) \ge \inf A_E.
      \end{align*}
      Por outro lado, dado $\nu(F) \in A_E$,
      \begin{align*}
        \nu(F) = \nu^+(F) - \nu^-(F) \ge -\nu^-(F) \ge -\nu^-(E).
      \end{align*}
    Desse modo, $\inf A_E \ge - \nu^-(E)$, ou seja $\nu^-(E) = -\inf A_E$, 
    como queríamos.
    \end{proof}
  \item $\mparth \nu (E) = \sup \set{\sum\limits_{j=1}^n \mparth{\nu(E_j)}, 
    \text{ em que $\set{E_j}_{j=1}^n$ é uma coleção finita disjunta com } X= \bigcup\limits_{j = 1}^n E_j}$. 
    \begin{proof}[\textbf{Dem:}]a 
    \end{proof}
\end{enumerate}}

%ex 4
\nextexe{seis}
\textbf{Sejam $\mu_1, \mu_2$ e $\mu_3$ medidas positivas num espaço mensurável $(X, \mcal F)$. 
Mostre que $\mu_1 \ll \mu_2$ e que $\mu_1 \ll \mu_2$ e $\mu_2 \ll \mu_3$
implicam $\mu_1 \ll \mu_3$. Dê um exemplo para mostrar que $\mu_1 \ll \mu_2$ 
não implica $\mu_2 \ll \mu_1$.}
\begin{proof}[\textbf{Dem:}]
  
\end{proof}
